\section*{SUMMARY}

Rhythmic behaviors, like swimming, rely on Central Pattern Generators (CPGs) that must be both robust and resilient to perturbation and tunable by descending commands.
Early work established that some CPGs act as ``network oscillators'' -- {\bf collectively generating emergent rhythm} through synaptic interactions rather than due to {\bf dedicated} cellular pacemakers alone \cite{Getting1983delayed, Getting1989}.
While pacemaker-driven and {\bf [bistable] half-center oscillators??} have been extensively studied, the mechanisms enabling \textit{smooth and broad frequency control} in network oscillators remain less well characterized.

We address this question through computational modeling of the swim CPG of the nudibranch \textit{Dendronotus iris}, a four-cell circuit controlled by Si1 command neurons.
We identify contralateral excitatory-inhibitory (E/I) modules as the core {\em rhythmogenic} engine {\bf symmetrically built-in the swim CPG}.

These modules operate in an ``almost-oscillatory'' {\bf [love it or on the oscillatory edge]} regime where rhythm emerges through network-level hysteresis: slow synaptic accumulation creates the history-dependence in post-synaptic neurons giving rise to  bistability and oscillations in E/I modules and the whole CPG. 

Progressive network assembly reveals that E/I modules, combined/enhanced with two half-center-oscillators between left/right constituent interneurons [Si2s and Si3s] bonded by reciprocal inhibition, constitute the circuit capable for emergent and resilient rhythm generation.

Its theoretical/mathematical model reproduces experimental perturbation responses without parameter fitting, confirming its distributed robustness or behavioral stability.
We further demonstrate that presynaptic modulation of transmitter release rate enables smooth 
network frequency control, whereas postsynaptic conductance modulation yields threshold-like transitions.
These findings provide a mechanistic account of how network oscillators achieve the {\bf graded}frequency control long observed in motor systems.
